% !TEX root =  mosaic.tex
\section{Evaluation}
\label{sec:eval}

We evaluate \name across three generations of NVIDIA GPUs whose hardware and software platforms are summarized in \autoref{tab:platforms}. The architectural parameters required by the cost model, including memory capacities, execution limits, and the profiled $L$ and $c$ constants described in \autoref{sec:gpu-target-profiling} and used by \autoref{sec:cost}, are gathered once per GPU through the target description of \autoref{sec:overview} and reused across all kernels.

\begin{table}[htbp]
  \centering
  \caption{Experimental platforms used in the evaluation.}
  \label{tab:platforms}
  \begin{tabular}{lcc}
    \toprule
    GPU       & Architecture & Memory \\
    \midrule
    RTX 4090  & Ada          & \SI{48}{GiB} \\
    A100      & Ampere       & \SI{80}{GiB} \\
    V100      & Volta        & \SI{32}{GiB} \\
    \midrule
    \multicolumn{3}{l}{\textit{Common software environment (all platforms)}} \\
    \addlinespace[2pt]
    OS            & \multicolumn{2}{l}{Ubuntu 22.04} \\
    CUDA Toolkit  & \multicolumn{2}{l}{CUDA 12.2} \\
    Compiler      & \multicolumn{2}{l}{NVCC release 12.2, V12.2.91} \\
    \bottomrule
  \end{tabular}
\end{table}

\name consumes schedule metadata from PPCG~0.09~\cite{ppcg-rep} for affine kernels and TVM~0.22.dev0~\cite{tvm-rep} for tensor operators and workloads. We compare against PPCG's fixed heuristic and its state-of-the-art EATSS analytical selector~\cite{eatss}, and the TVM MetaSchedule autotuner~\cite{metaschedule}. To isolate execution space selection, \name operates on the same transformed programs as each baseline; only tile sizes, block dimensions, and shared memory placement differ.

\subsection{Performance on PolyBench Kernels}
\label{subsec:polybench}

% polybench table
\begin{table*}[htbp]
\centering
\caption{Execution time, solver time, and speedups on PolyBench benchmarks.}
\label{tab:benchmark}
\scriptsize
\setlength{\tabcolsep}{2.4pt}
\renewcommand{\arraystretch}{0.86}
\adjustbox{max width=\textwidth}{
\begin{tabular}{lccccccccccccccc}
\toprule
\multirow{3}{*}{Benchmark} & \multicolumn{5}{c}{RTX 4090} & \multicolumn{5}{c}{A100} & \multicolumn{5}{c}{V100} \\
\cmidrule(lr){2-6} \cmidrule(lr){7-11} \cmidrule(lr){12-16}
& \multicolumn{2}{c}{Mosaic} & \multicolumn{1}{c}{PPCG} & \multicolumn{2}{c}{EATSS}
& \multicolumn{2}{c}{Mosaic} & \multicolumn{1}{c}{PPCG} & \multicolumn{2}{c}{EATSS}
& \multicolumn{2}{c}{Mosaic} & \multicolumn{1}{c}{PPCG} & \multicolumn{2}{c}{EATSS} \\
\cmidrule(lr){2-3} \cmidrule(lr){4-4} \cmidrule(lr){5-6}
\cmidrule(lr){7-8} \cmidrule(lr){9-9} \cmidrule(lr){10-11}
\cmidrule(lr){12-13} \cmidrule(lr){14-14} \cmidrule(lr){15-16}
& Exec. (ms) & Solver (s) & Exec. Spd. & Exec. Spd. & Solver Spd.
& Exec. (ms) & Solver (s) & Exec. Spd. & Exec. Spd. & Solver Spd.
& Exec. (ms) & Solver (s) & Exec. Spd. & Exec. Spd. & Solver Spd. \\
\midrule
correlation & 16.578 & 0.216 & 1.07 & 1.56 & 6.42 & 10.383 & 0.362 & 1.05 & 4.55 & 4.86 & 10.366 & 0.464 & 1.05 & 3.52 & 3.09 \\
covariance & 17.247 & 0.213 & 2.11 & 2.54 & 6.17 & 11.520 & 0.143 & 3.62 & 4.35 & 11.5 & 41.364 & 0.146 & 3.30 & 1.72 & 9.72 \\
gemm & 22.704 & 0.242 & 1.39 & 1.08 & 5.71 & 7.8790 & 0.245 & 1.23 & 9.52 & 6.45 & 14.269 & 0.243 & 2.01 & 4.24 & 6.10 \\
gemver & 0.7540 & 0.050 & 1.01 & 1.03 & 34.5 & 1.1800 & 0.050 & 1.04 & 1.17 & 35.3 & 1.1980 & 0.050 & 1.04 & 1.50 & 34.4 \\
gesummv & 0.5570 & 0.006 & 1.00 & 1.19 & 284 & 0.5410 & 0.011 & 1.71 & 0.90 & 141 & 0.5430 & 0.011 & 1.50 & 1.02 & 158 \\
symm & 22.350 & 0.030 & 2.27 & 2.19 & 52.1 & 18.954 & 0.162 & 6.05 & 2.88 & 9.37 & 63.452 & 0.164 & 6.02 & 1.64 & 9.40 \\
syr2k & 58.963 & 3.324 & 1.36 & 1.35 & 0.45 & 49.926 & 0.127 & 1.65 & 1.91 & 11.9 & 101.02 & 0.080 & 2.70 & 1.30 & 19.8 \\
syrk & 19.681 & 0.069 & 1.12 & 1.27 & 18.6 & 8.3080 & 0.064 & 5.58 & 5.56 & 13.4 & 41.736 & 0.061 & 3.67 & 1.77 & 21.4 \\
trmm & 136.68 & 0.099 & 1.03 & 2.51 & 17.4 & 139.91 & 0.083 & 0.93 & 2.69 & 18.8 & 331.15 & 0.084 & 1.11 & 1.90 & 20.8 \\
2mm & 22.380 & 0.266 & 1.31 & 1.23 & 5.27 & 7.7950 & 0.260 & 1.50 & 4.07 & 7.11 & 24.181 & 0.265 & 1.32 & 1.79 & 5.60 \\
3mm & 32.569 & 0.431 & 1.09 & 1.31 & 3.37 & 15.573 & 0.442 & 1.26 & 2.38 & 3.61 & 41.091 & 0.426 & 1.23 & 1.26 & 3.54 \\
atax & 0.1920 & 0.077 & 1.21 & 1.14 & 22.2 & 0.3730 & 0.061 & 1.51 & 0.83 & 31.7 & 0.3720 & 0.061 & 1.54 & 1.46 & 28.6 \\
bicg & 0.1920 & 0.021 & 1.16 & 1.09 & 84.7 & 0.3740 & 0.066 & 1.49 & 0.84 & 25.5 & 0.3790 & 0.022 & 1.44 & 1.30 & 81.2 \\
doitgen & 702.24 & 0.001 & 1.01 & -- & -- & 828.13 & 0.043 & 1.17 & -- & -- & 689.50 & 0.022 & 1.61 & -- & -- \\
mvt & 0.5290 & 0.019 & 1.00 & 1.92 & 84.5 & 0.9420 & 0.010 & 0.98 & 1.21 & 181 & 0.9520 & 0.009 & 0.98 & 1.16 & 186 \\
cholesky & 140.19 & 0.090 & 1.14 & 1.53 & 18.6 & 191.94 & 0.128 & 1.25 & 1.81 & 13.1 & 356.29 & 0.116 & 1.07 & 1.65 & -- \\
gramschmidt & 685.11 & 0.077 & 1.14 & 1.26 & 19.9 & 1136.7 & 0.092 & 1.11 & 1.09 & 17.7 & 1398.8 & 0.144 & 1.11 & 1.13 & 10.5 \\
lu & 330.18 & 0.215 & 1.31 & 1.00 & 7.98 & 321.98 & 0.222 & 1.90 & 1.00 & 6.97 & 638.70 & 0.221 & 1.44 & 1.04 & 7.59 \\
trisolv & 9.9610 & 0.019 & 1.00 & 1.00 & 67.2 & 16.099 & 0.019 & 1.00 & 1.00 & 79.3 & 14.542 & 0.019 & 1.00 & 1.00 & 66.1 \\
adi & 4090.7 & 0.580 & 1.51 & 1.00 & 2.68 & 3406.3 & 0.574 & 1.40 & 1.01 & 3.28 & 3126.2 & 0.572 & 1.33 & 1.01 & 2.77 \\
fdtd-2d & 452.12 & 0.652 & 1.03 & 1.57 & 2.34 & 301.55 & 0.663 & 1.06 & 2.30 & 2.75 & 585.86 & 0.563 & 1.02 & 1.59 & 2.87 \\
heat-3d & 881.13 & 1.129 & 1.58 & -- & -- & 662.45 & 1.158 & 1.10 & -- & -- & 172.49 & 1.105 & 1.66 & -- & -- \\
jacobi-1d & 2.6740 & 0.078 & 1.00 & -- & -- & 4.1800 & 0.079 & 1.02 & -- & -- & 3.3060 & 0.079 & 1.06 & -- & -- \\
jacobi-2d & 266.16 & 0.312 & 1.00 & 1.17 & 5.14 & 178.16 & 0.311 & 1.01 & 1.35 & 5.56 & 422.13 & 0.312 & 1.00 & 2.02 & 4.98 \\
seidel-2d & 757.94 & 0.270 & 1.27 & 1.51 & 5.63 & 935.38 & 0.241 & 1.01 & 1.41 & 7.51 & 1726.5 & 0.244 & 1.19 & 1.79 & 6.53 \\
\midrule
\textbf{geomean} & 346.79 & 0.34 & \textbf{1.21} & \textbf{1.37} & \textbf{12.6} & 330.26 & 0.22 & \textbf{1.44} & \textbf{1.89} & \textbf{13.6} & 392.25 & 0.22 & \textbf{1.49} & \textbf{1.55} & \textbf{13.9} \\
\bottomrule
\end{tabular}
}
\end{table*}

We first evaluate \name on twenty-five PolyBench kernels~\cite{polybench} spanning reductions, contractions, and stencils, comparing against the PPCG heuristic and EATSS on all three GPUs. \autoref{tab:benchmark} reports, for each platform, \name's execution and solver times together with its speedups over PPCG and EATSS and the reduction in selection overhead relative to EATSS. Speedup means are geometric while solver-time means are arithmetic. \name consistently outperforms both baselines across all three architectures. Over PPCG, it achieves mean speedups of $1.21\times$, $1.44\times$, and $1.49\times$ on the RTX 4090, A100, and V100, respectively; over EATSS, it achieves $1.37\times$, $1.89\times$, and $1.55\times$. Gains are largest on contraction-heavy kernels, where the \autoref{sec:bg-challenges} loop is most pronounced: on the A100, \name improves \texttt{symm} by $6.05\times$ and \texttt{syrk} by $5.58\times$ over PPCG. The advantage holds as register and shared memory budgets vary across GPU generations.

\name is occasionally within \(5\%\) of EATSS on the A100 for the small memory-bound kernels \texttt{atax}, \texttt{bicg}, and \texttt{gesummv}, where the analytical and searched optima nearly coincide. The \texttt{trmm} kernel on A100 is marginally below PPCG (\(0.93\times\)) and \texttt{trisolv} is essentially tied with both baselines (\(\approx1.00\times\)) on every platform, reflecting limited headroom for execution space tuning. EATSS fails to produce configurations for \texttt{doitgen}, \texttt{heat-3d}, and \texttt{jacobi-1d}, shown as ``--'' in \autoref{tab:benchmark}.

\subsection{Performance on Deep Learning Operators}
\label{subsec:operators}

We next evaluate \name on five deep learning operators. CNN, LSTM, and RNN come from the PolyNN suite~\cite{polynn}; FNN and Embedding are extracted from the end-to-end network workloads. \autoref{fig:network-operators} reports the speedup of \name over PPCG. \name delivers large gains on CNN (\(5.73\times\)) and Embedding (\(3.85\times\)) and remains within \(4\%\) of PPCG on LSTM, RNN, and FNN, yielding a geometric-mean speedup of \(1.88\times\).

\begin{figure}[!tbp]
\centering
\begin{tikzpicture}
\begin{groupplot}[
    group style={
        group size=1 by 2,
        vertical sep=3pt,
        xticklabels at=edge bottom,
    },
    ybar,
    /pgf/bar width=5pt,        
    width=0.9\columnwidth,
    symbolic x coords={CNN,LSTM,RNN,FNN,Embed},
    xtick=data,
    enlarge x limits=0.13,
    ymajorgrids=true,
    grid style=dashed,
    xticklabel style={font=\scriptsize, text depth=0pt},
    yticklabel style={font=\scriptsize},
    nodes near coords,
    every node near coord/.append style={font=\tiny, rotate=90, anchor=west},
    point meta=explicit symbolic,
]

\nextgroupplot[
    height=0.24\columnwidth,
    ymin=3, ymax=7.2,
    ytick={4.00,6.00}, yticklabels={4.00,6.00},
    legend style={
        at={(0.5,2.5)}, anchor=north,
        legend columns=2, font=\scriptsize,
        /tikz/every even column/.append style={column sep=0.15cm},
    },
]
\addplot[fill=color1, draw=black] coordinates {
    (CNN,1)[] (LSTM,1)[] (RNN,1)[] (FNN,1)[] (Embed,1)[]
};
\addplot[fill=color4, draw=black] coordinates {
    (CNN,5.734)[$\mathbf{5.734}\mathbf{\times}$] (LSTM,1.015)[] (RNN,1.035)[] (FNN,1.002)[] (Embed,3.853)[$\mathbf{3.853}\mathbf{\times}$]
};
\legend{PPCG, Mosaic}

\nextgroupplot[
    height=0.34\columnwidth,
    ymin=0.985, ymax=1.06,
    ytick={1.00,1.02,1.04}, yticklabels={1.00,1.02,1.04},
]
\addplot[fill=color1, draw=black] coordinates {
    (CNN,1)[549.69ms] (LSTM,1)[1010.2ms] (RNN,1)[1452.4ms] (FNN,1)[1.6980ms] (Embed,1)[544.26ms]
};
\addplot[fill=color4, draw=black] coordinates {
    (CNN,5.734)[] (LSTM,1.015)[$\mathbf{1.015}\mathbf{\times}$] (RNN,1.035)[$\mathbf{1.035}\mathbf{\times}$] (FNN,1.002)[$\mathbf{1.002}\mathbf{\times}$] (Embed,3.853)[]
};
\end{groupplot}


\coordinate (bl) at ($(group c1r1.south west)!0.5!(group c1r2.north west)$);
\coordinate (br) at ($(group c1r1.south east)!0.5!(group c1r2.north east)$);
\foreach \p in {bl,br}{
    \draw[white,line width=2pt] ($(\p)+(0,-3pt)$)--($(\p)+(0,3pt)$);
    \draw[line width=.5pt] ($(\p)+(-3pt,-3pt)$)--($(\p)+(3pt,0pt)$);
    \draw[line width=.5pt] ($(\p)+(-3pt,0pt)$)--($(\p)+(3pt,3pt)$);
}


\node[rotate=90, anchor=south, font=\scriptsize]
    at ([xshift=-20pt]$(group c1r1.west)!0.5!(group c1r2.west)$)
    {Execution time speedup};
\end{tikzpicture}
\caption{Performance of deep learning operators.}
\Description{Bar chart comparing PPCG and \name speedups on CNN, LSTM, RNN, FNN, and embedding operators.}
\label{fig:network-operators}
\end{figure}

These operators cover different execution regimes. CNN and FNN are contraction/GEMM-like and therefore compute-intensive, but FNN is small and regular enough that the PPCG default already lands close to the configuration selected by \name. LSTM, in contrast, is a representative memory-intensive recurrent operator; its performance is dominated more by data movement than by arithmetic intensity, leaving less room for the register-blocking trade-offs captured by the $\phi$ ratio. Embedding is also data-movement dominated, but its default PPCG mapping leaves substantial opportunity for \name's placement and occupancy model to improve the execution configuration.

To compare \name against search-based autotuning, we integrate it into TVM~\cite{tvm} and compare with MetaSchedule~\cite{metaschedule} across trial budgets. \autoref{fig:network-operator-curves} plots execution time and selection cost for six operators. \name is flat because it solves once, whereas MetaSchedule needs hundreds to thousands of trials to reach comparable execution time and pays search cost that grows with the budget. \name therefore attains autotuner-level performance while reducing selection overhead by \(2\)--\(4\) orders of magnitude.

\begin{figure}[!tbp]
\centering
\definecolor{curveMosaic}{HTML}{CC0000}
\definecolor{curveTVM}{HTML}{FF6666}
\definecolor{curveSolverMosaic}{HTML}{0066CC}
\definecolor{curveSolverTVM}{HTML}{4D94E6}
\begin{tikzpicture}
\begin{scope}[xshift=0.1\columnwidth, yshift=4.55cm]
\begin{axis}[
    width=0.45\columnwidth, height=2.35cm,
    yticklabel style={color=curveMosaic, font=\scriptsize},
    ymode=log,
    ymin=1, ymax=100,
    ytick={1,10,100},
    xmode=log,
    xmin=0.8, xmax=1200,
    xtick={1,10,100,1000},
    xticklabels={1,$10^1$,$10^2$,$10^3$},
    xticklabel style={font=\scriptsize},
    grid=major,
    grid style={line width=0.3pt, draw=gray!25},
    axis y line*=left,
    axis x line*=bottom,
    title={CNN},
    title style={font=\footnotesize, yshift=-0.2cm},
]
\addplot[color=curveMosaic, line width=0.55pt, solid, mark=*, mark size=1.2pt] coordinates {
    (1,4.2312) (10,4.2312) (100,4.2312) (200,4.2312) (500,4.2312) (1000,4.2312)
};
\addplot[color=curveTVM, line width=0.55pt, solid, mark=square*, mark size=1.2pt] coordinates {
    (1,51.732) (10,60.3056) (100,6.7591) (200,6.8391) (500,7.4544) (1000,5.8638)
};
\end{axis}
\begin{axis}[
    width=0.45\columnwidth, height=2.35cm,
    axis y line*=right,
    axis x line=none,
    yticklabel style={color=curveSolverMosaic, opacity=0},
    ymode=log,
    ymin=0.1, ymax=1000,
    ytick={0.1,10,1000},
    xmode=log,
    xmin=0.8, xmax=1200,
]
\addplot[color=curveSolverMosaic, line width=0.55pt, dashed, mark=*, mark size=1.2pt] coordinates {
    (1,2.9676) (10,2.9676) (100,2.9676) (200,2.9676) (500,2.9676) (1000,2.9676)
};
\addplot[color=curveSolverTVM, line width=0.55pt, dashed, mark=square*, mark size=1.2pt] coordinates {
    (1,17.2) (10,42.3) (100,737.7) (200,939.3) (500,924.3) (1000,1078)
};
\end{axis}
\end{scope}

\begin{scope}[xshift=0.55\columnwidth, yshift=4.55cm]
\begin{axis}[
    width=0.45\columnwidth, height=2.35cm,
    yticklabel style={color=curveMosaic, font=\scriptsize},
    ymode=log,
    ymin=0.1, ymax=100,
    ytick={0.1,1,10},
    xmode=log,
    xmin=0.8, xmax=1200,
    xtick={1,10,100,1000},
    xticklabels={1,$10^1$,$10^2$,$10^3$},
    xticklabel style={font=\scriptsize},
    grid=major,
    grid style={line width=0.3pt, draw=gray!25},
    axis y line*=left,
    axis x line*=bottom,
    title={LSTM},
    title style={font=\footnotesize, yshift=-0.2cm},
]
\addplot[color=curveMosaic, line width=0.55pt, solid, mark=*, mark size=1.2pt] coordinates {
    (1,0.2332) (10,0.2332) (100,0.2332) (200,0.2332) (500,0.2332) (1000,0.2332)
};
\addplot[color=curveTVM, line width=0.55pt, solid, mark=square*, mark size=1.2pt] coordinates {
    (1,21.9229) (10,1.5843) (100,0.2814) (200,0.3319) (500,0.2185) (1000,0.2185)
};
\end{axis}
\begin{axis}[
    width=0.45\columnwidth, height=2.35cm,
    axis y line*=right,
    axis x line=none,
    yticklabel style={color=curveSolverMosaic, font=\scriptsize},
    ymode=log,
    ymin=0.01, ymax=2000,
    ytick={0.01,10,1000},
    xmode=log,
    xmin=0.8, xmax=1200,
]
\addplot[color=curveSolverMosaic, line width=0.55pt, dashed, mark=*, mark size=1.2pt] coordinates {
    (1,0.0741) (10,0.0741) (100,0.0741) (200,0.0741) (500,0.0741) (1000,0.0741)
};
\addplot[color=curveSolverTVM, line width=0.55pt, dashed, mark=square*, mark size=1.2pt] coordinates {
    (1,16) (10,22.47) (100,131.84) (200,267.54) (500,660.27) (1000,1326.86)
};
\end{axis}
\end{scope}

\begin{scope}[xshift=0.1\columnwidth, yshift=3cm]
\begin{axis}[
    width=0.45\columnwidth, height=2.35cm,
    yticklabel style={color=curveMosaic, font=\scriptsize},
    ymode=log,
    ymin=10, ymax=3000,
    ytick={10,100,1000},
    xmode=log,
    xmin=0.8, xmax=1200,
    xtick={1,10,100,1000},
    xticklabels={1,$10^1$,$10^2$,$10^3$},
    xticklabel style={font=\scriptsize},
    grid=major,
    grid style={line width=0.3pt, draw=gray!25},
    axis y line*=left,
    axis x line*=bottom,
    title={RNN},
    title style={font=\footnotesize, yshift=-0.2cm},
]
\addplot[color=curveMosaic, line width=0.55pt, solid, mark=*, mark size=1.2pt] coordinates {
    (1,20.4271) (10,20.4271) (100,20.4271) (200,20.4271) (500,20.4271) (1000,20.4271)
};
\addplot[color=curveTVM, line width=0.55pt, solid, mark=square*, mark size=1.2pt] coordinates {
    (1,2252.1594) (10,81.9091) (100,21.8536) (200,20.0876) (500,17.8522) (1000,17.7766)
};
\end{axis}
\begin{axis}[
    width=0.45\columnwidth, height=2.35cm,
    axis y line*=right,
    axis x line=none,
    yticklabel style={color=curveSolverMosaic, opacity=0},
    ymode=log,
    ymin=0.1, ymax=2000,
    ytick={0.1,10,1000},
    xmode=log,
    xmin=0.8, xmax=1200,
]
\addplot[color=curveSolverMosaic, line width=0.55pt, dashed, mark=*, mark size=1.2pt] coordinates {
    (1,0.716) (10,0.716) (100,0.716) (200,0.716) (500,0.716) (1000,0.716)
};
\addplot[color=curveSolverTVM, line width=0.55pt, dashed, mark=square*, mark size=1.2pt] coordinates {
    (1,23.44) (10,219.66) (100,823.3) (200,895.68) (500,1007.83) (1000,1012.12)
};
\end{axis}
\end{scope}

\begin{scope}[xshift=0.55\columnwidth, yshift=3cm]
\begin{axis}[
    width=0.45\columnwidth, height=2.35cm,
    yticklabel style={color=curveMosaic, font=\scriptsize},
    ymode=log,
    ymin=0.05, ymax=2,
    ytick={0.05,0.2,1},
    xmode=log,
    xmin=0.8, xmax=1200,
    xtick={1,10,100,1000},
    xticklabels={1,$10^1$,$10^2$,$10^3$},
    xticklabel style={font=\scriptsize},
    grid=major,
    grid style={line width=0.3pt, draw=gray!25},
    axis y line*=left,
    axis x line*=bottom,
    title={FNN},
    title style={font=\footnotesize, yshift=-0.2cm},
]
\addplot[color=curveMosaic, line width=0.55pt, solid, mark=*, mark size=1.2pt] coordinates {
    (1,0.0738) (10,0.0738) (100,0.0738) (200,0.0738) (500,0.0738) (1000,0.0738)
};
\addplot[color=curveTVM, line width=0.55pt, solid, mark=square*, mark size=1.2pt] coordinates {
    (1,0.9439) (10,1.2417) (100,0.1618) (200,0.0944) (500,0.081) (1000,0.0741)
};
\end{axis}
\begin{axis}[
    width=0.45\columnwidth, height=2.35cm,
    axis y line*=right,
    axis x line=none,
    yticklabel style={color=curveSolverMosaic, font=\scriptsize},
    ymode=log,
    ymin=0.01, ymax=1000,
    ytick={0.01,1,100},
    xmode=log,
    xmin=0.8, xmax=1200,
]
\addplot[color=curveSolverMosaic, line width=0.55pt, dashed, mark=*, mark size=1.2pt] coordinates {
    (1,0.1194) (10,0.1194) (100,0.1194) (200,0.1194) (500,0.1194) (1000,0.1194)
};
\addplot[color=curveSolverTVM, line width=0.55pt, dashed, mark=square*, mark size=1.2pt] coordinates {
    (1,9.8) (10,15.2) (100,78.2) (200,122.8) (500,282.5) (1000,561.9)
};
\end{axis}
\end{scope}

\begin{scope}[xshift=0.1\columnwidth, yshift=1.4cm]
\begin{axis}[
    width=0.45\columnwidth, height=2.35cm,
    yticklabel style={color=curveMosaic, font=\scriptsize},
    ymode=log,
    ymin=10, ymax=500,
    ytick={10,50,200},
    xmode=log,
    xmin=0.8, xmax=1200,
    xtick={1,10,100,1000},
    xticklabels={1,$10^1$,$10^2$,$10^3$},
    xticklabel style={font=\scriptsize},
    grid=major,
    grid style={line width=0.3pt, draw=gray!25},
    axis y line*=left,
    axis x line*=bottom,
    title={Embedding},
    title style={font=\footnotesize, yshift=-0.2cm},
]
\addplot[color=curveMosaic, line width=0.55pt, solid, mark=*, mark size=1.2pt] coordinates {
    (1,12.9526) (10,12.9526) (100,12.9526) (200,12.9526) (500,12.9526) (1000,12.9526)
};
\addplot[color=curveTVM, line width=0.55pt, solid, mark=square*, mark size=1.2pt] coordinates {
    (1,362.8838) (10,85.9239) (100,38.2759) (200,32.3857) (500,29.9019) (1000,25.491)
};
\end{axis}
\begin{axis}[
    width=0.45\columnwidth, height=2.35cm,
    axis y line*=right,
    axis x line=none,
    yticklabel style={color=curveSolverMosaic, opacity=0},
    ymode=log,
    ymin=0.01, ymax=3000,
    ytick={0.01,10,1000},
    xmode=log,
    xmin=0.8, xmax=1200,
]
\addplot[color=curveSolverMosaic, line width=0.55pt, dashed, mark=*, mark size=1.2pt] coordinates {
    (1,0.3567) (10,0.3567) (100,0.3567) (200,0.3567) (500,0.3567) (1000,0.3567)
};
\addplot[color=curveSolverTVM, line width=0.55pt, dashed, mark=square*, mark size=1.2pt] coordinates {
    (1,16.71) (10,108.89) (100,610.7) (200,1041.74) (500,2432.51) (1000,2432.51)
};
\end{axis}
\end{scope}

\begin{scope}[xshift=0.55\columnwidth, yshift=1.4cm]
\begin{axis}[
    width=0.45\columnwidth, height=2.35cm,
    yticklabel style={color=curveMosaic, font=\scriptsize},
    ymode=log,
    ymin=1, ymax=50,
    ytick={1,5,20},
    xmode=log,
    xmin=0.8, xmax=1200,
    xtick={1,10,100,1000},
    xticklabels={1,$10^1$,$10^2$,$10^3$},
    xticklabel style={font=\scriptsize},
    grid=major,
    grid style={line width=0.3pt, draw=gray!25},
    axis y line*=left,
    axis x line*=bottom,
    title={Matmul},
    title style={font=\footnotesize, yshift=-0.2cm},
]
\addplot[color=curveMosaic, line width=0.55pt, solid, mark=*, mark size=1.2pt] coordinates {
    (1,4.4347) (10,4.4347) (100,4.4347) (200,4.4347) (500,4.4347) (1000,4.4347)
};
\addplot[color=curveTVM, line width=0.55pt, solid, mark=square*, mark size=1.2pt] coordinates {
    (1,26.5995) (10,20.8549) (100,7.149) (200,4.2035) (500,4.5826) (1000,4.1599)
};
\end{axis}
\begin{axis}[
    width=0.45\columnwidth, height=2.35cm,
    axis y line*=right,
    axis x line=none,
    yticklabel style={color=curveSolverMosaic, font=\scriptsize},
    ymode=log,
    ymin=0.01, ymax=2000,
    ytick={0.01,10,1000},
    xmode=log,
    xmin=0.8, xmax=1200,
]
\addplot[color=curveSolverMosaic, line width=0.55pt, dashed, mark=*, mark size=1.2pt] coordinates {
    (1,0.21) (10,0.21) (100,0.21) (200,0.21) (500,0.21) (1000,0.21)
};
\addplot[color=curveSolverTVM, line width=0.55pt, dashed, mark=square*, mark size=1.2pt] coordinates {
    (1,11.06) (10,88.17) (100,360.63) (200,601.6) (500,866.03) (1000,1671.5)
};
\end{axis}
\end{scope}

\node[rotate=90, color=curveMosaic, font=\scriptsize] at (0\columnwidth, 3.38cm) {Execution time (ms)};
\node[rotate=90, color=curveSolverMosaic, font=\scriptsize] at (0.91\columnwidth, 3.38cm) {Search/solver (s)};
\node[font=\scriptsize] at (0.47\columnwidth, 0.80cm) {Number of Trials};
\node at (0.48\columnwidth, 0.42cm) {
    \begin{tikzpicture}
        % Mosaic exec.
        \draw[color=curveMosaic, line width=0.7pt, solid] (0,0) -- (0.28,0);
        \filldraw[fill=curveMosaic, draw=white, line width=0.3pt]
            (0.16,0) circle (0.045);
        \node[right, font=\tiny] at (0.34,0) {Mosaic exec.};

        % TVM exec.
        \draw[color=curveTVM, line width=0.7pt, solid] (1.57,0) -- (1.85,0);
        \filldraw[fill=curveTVM, draw=white, line width=0.3pt]
            (1.68,-0.05) rectangle (1.78,0.05);
        \node[right, font=\tiny] at (1.89,0) {TVM exec.};

        % Mosaic solver
        \draw[color=curveSolverMosaic, line width=0.7pt, dashed] (3.00,0) -- (3.28,0);
        \filldraw[fill=curveSolverMosaic, draw=white, line width=0.3pt]
            (3.16,0) circle (0.045);
        \node[right, font=\tiny] at (3.34,0) {Mosaic solver};

        % TVM search
        \draw[color=curveSolverTVM, line width=0.7pt, dashed] (4.75,0) -- (5.03,0);
        \filldraw[fill=curveSolverTVM, draw=white, line width=0.3pt]
            (4.86,-0.05) rectangle (4.96,0.05);
        \node[right, font=\tiny] at (5.09,0) {TVM search};
    \end{tikzpicture}
};

\end{tikzpicture}
\caption{Execution time and search/solver overhead for six operators, comparing with TVM MetaSchedule. Solid lines show execution time; dashed lines show selection cost.}
\Description{Six curve plots for CNN, LSTM, RNN, FNN, embedding, and matmul comparing \name and tvm SSSRRSRS execution time and search or solver overhead.}
\label{fig:network-operator-curves}
\end{figure}

\subsection{Performance on End-to-End Workloads}
\label{subsec:e2e}

We compose the operators into three end-to-end workloads spanning the major network families: AlexNet~\cite{alexnet} (CNN), Seq2Seq~\cite{lstm} (RNN), and BertTiny~\cite{bert} (Transformer). \autoref{fig:e2e-fourbar} reports the per-network speedup of \name over both PPCG and TVM, with the baseline execution times shown on the left of the corresponding bars. Over PPCG, \name achieves $1.66\times$, $9.23\times$, and $4.90\times$ on AlexNet, Seq2Seq, and BertTiny, respectively, for a mean speedup of $4.22\times$. Over TVM MetaSchedule, \name achieves $1.16\times$, $1.05\times$, and $1.13\times$, for a mean of $1.11\times$. These results demonstrate that \name slightly exceeds a state-of-the-art autotuner on full workloads without performing any architecture-space search, while delivering large improvements over the heuristic configurations of PPCG. The end-to-end gains over PPCG are more pronounced than the per-operator gains of \autoref{subsec:operators}, as consistently good execution space choices across consecutive layers compound into a larger overall improvement.

\begin{figure}[htbp]
\centering
\begin{tikzpicture}
\begin{groupplot}[
    group style={
        group size=1 by 2,
        vertical sep=3pt,
        xticklabels at=edge bottom,
    },
    ybar,
    /pgf/bar width=6pt,
    width=0.9\columnwidth,
    symbolic x coords={AlexNet,Seq2Seq,BertTiny},
    xtick=data,
    enlarge x limits=0.18,
    ymajorgrids=true,
    grid style=dashed,
    xticklabel style={font=\scriptsize},
    yticklabel style={font=\scriptsize},
    nodes near coords,
    every node near coord/.append style={font=\tiny, rotate=90, anchor=west},
    point meta=explicit symbolic,
]

\nextgroupplot[
    height=0.34\columnwidth,
    ymin=1.4, ymax=14,
    ytick={2,4,6,8,10,12}, yticklabels={2,4,6,8,10,12},
    legend style={at={(0.5,1.18)}, anchor=south, legend columns=4, font=\scriptsize,
        /tikz/every even column/.append style={column sep=0.12cm}},
]
\addplot[fill=color1,draw=black] coordinates {(AlexNet,1)[] (Seq2Seq,1)[] (BertTiny,1)[]};
\addplot[fill=color4,draw=black] coordinates {(AlexNet,1.660)[\textbf{1.660$\times$}] (Seq2Seq,9.227)[\textbf{9.227$\times$}] (BertTiny,4.900)[\textbf{4.900$\times$}]};
\addplot[fill=color2,draw=black] coordinates {(AlexNet,1)[] (Seq2Seq,1)[] (BertTiny,1)[]};
\addplot[fill=color4,draw=black] coordinates {(AlexNet,1.159)[] (Seq2Seq,1.052)[] (BertTiny,1.127)[]};
\legend{PPCG, \name/PPCG, TVM, \name/TVM}

\nextgroupplot[
    height=0.30\columnwidth,
    ymin=0.9, ymax=1.22,
    ytick={1.00,1.10,1.20}, yticklabels={1.00,1.10,1.20},
]
\addplot[fill=color1,draw=black] coordinates {(AlexNet,1)[850.869ms] (Seq2Seq,1)[102720ms] (BertTiny,1)[226.710ms]};
\addplot[fill=color4,draw=black] coordinates {(AlexNet,1.660)[] (Seq2Seq,9.227)[] (BertTiny,4.900)[]};
\addplot[fill=color2,draw=black] coordinates {(AlexNet,1)[26.2720ms] (Seq2Seq,1)[614.433ms] (BertTiny,1)[5.40400ms]};
\addplot[fill=color4,draw=black] coordinates {(AlexNet,1.159)[\textbf{1.159$\times$}] (Seq2Seq,1.052)[\textbf{1.052$\times$}] (BertTiny,1.127)[\textbf{1.127$\times$}]};
\end{groupplot}
 

\coordinate (bl) at ($(group c1r1.south west)!0.5!(group c1r2.north west)$);
\coordinate (br) at ($(group c1r1.south east)!0.5!(group c1r2.north east)$);
\foreach \p in {bl,br}{
    \draw[white,line width=2pt] ($(\p)+(0,-3pt)$)--($(\p)+(0,3pt)$);
    \draw[line width=.5pt] ($(\p)+(-3pt,-3pt)$)--($(\p)+(3pt,0pt)$);
    \draw[line width=.5pt] ($(\p)+(-3pt,0pt)$)--($(\p)+(3pt,3pt)$);
}
\node[rotate=90, anchor=south, font=\scriptsize]
    at ([xshift=-20pt]$(group c1r1.west)!0.5!(group c1r2.west)$) {Execution time speedup};
\end{tikzpicture}
\caption{Performance of end-to-end workloads.}
\Description{Four bars per network: PPCG baseline, \name over PPCG, TVM baseline, and \name over TVM.}
\label{fig:e2e-fourbar}
\end{figure}

\subsection{Ablation Study}
\label{subsec:ablation}

We ablate the main modeling components on representative kernels and operators. FP-Wave uses the quantized wave objective with a fixed placement; RF-Wave adds the conservative register-feasibility filter; the full \name model further enumerates shared memory placement schemes. Using FP-Wave as the baseline, RF-Wave improves performance by a  mean of $1.37\times$, while full \name reaches $1.84\times$. The largest gain comes on CNN, where register filtering and placement selection together improve the selected configuration from $265.793$~ms to $95.851$~ms.

\begin{figure}[htbp]
\centering
\begin{tikzpicture}
\begin{axis}[
    ybar,
    /pgf/bar width=5pt,
    width=0.9\columnwidth,
    height=0.40\columnwidth,
    symbolic x coords={covariance,syrk,3mm,gemm,cnn},
    xtick=data,
    enlarge x limits=0.15,
    ymin=0.8,
    ymax=3.0,
    ytick={1.0,1.5,2.0,2.5},
    yticklabels={1.0,1.5,2.0,2.5},
    ymajorgrids=true,
    grid style=dashed,
    xticklabel style={font=\scriptsize},
    yticklabel style={font=\scriptsize},
    ylabel={Execution time speedup},
    ylabel style={font=\scriptsize},
    nodes near coords,
    every node near coord/.append style={font=\tiny, rotate=90, anchor=west},
    point meta=explicit symbolic,
    legend style={
        at={(0.5,1.10)}, anchor=south,
        legend columns=3, font=\scriptsize,
        /tikz/every even column/.append style={column sep=0.12cm}
    },
]
\addplot[fill=color1, draw=black] coordinates {
    (covariance,1)[31.512ms]
    (syrk,1)[28.324ms]
    (3mm,1)[50.674ms]
    (gemm,1)[36.994ms]
    (cnn,1)[265.793ms]
};
\addplot[fill=color3, draw=black] coordinates {
    (covariance,1.152)[\textbf{1.152$\times$}]
    (syrk,1.319)[\textbf{1.319$\times$}]
    (3mm,1.004)[\textbf{1.004$\times$}]
    (gemm,1.125)[\textbf{1.125$\times$}]
    (cnn,2.337)[\textbf{2.337$\times$}]
};
\addplot[fill=color4, draw=black] coordinates {
    (covariance,1.827)[\textbf{1.827$\times$}]
    (syrk,1.439)[\textbf{1.439$\times$}]
    (3mm,1.556)[\textbf{1.556$\times$}]
    (gemm,1.630)[\textbf{1.630$\times$}]
    (cnn,2.773)[\textbf{2.773$\times$}]
};
\legend{FP-Wave, RF-Wave, Mosaic}
\end{axis}
\end{tikzpicture}
\caption{Ablation study on \name's individual components.}
\Description{Grouped bar chart showing execution-time speedups of RF-Wave and Mosaic over the FP-Wave baseline on covariance, syrk, 3mm, gemm, and CNN.}
\label{fig:ablation}
\end{figure}

\subsection{Solver Overhead and Scalability}
\label{subsec:overhead}

Finally, we examine the cost of \name's analytical selection, reported in the solver-time columns of \autoref{tab:benchmark}. \name completes its two-stage solving in well under one second for the large majority of kernels and in at most a few seconds in the most demanding case (e.g., \texttt{syr2k} at $3.32$~s on the RTX 4090), since the feasible domain is bounded by the guard-free candidate sets of \autoref{sec:abstraction} and pruned by the conservative filters of \autoref{sec:constraints} before the MILP is ever invoked. Relative to EATSS, \name reduces selection overhead by mean factors of $12.6\times$, $13.6\times$, and $13.9\times$ on RTX 4090, A100, and V100, respectively, and the gap relative to the trial-based search of TVM spans several orders of magnitude (\autoref{fig:network-operator-curves}). As with the per-architecture profile, this selection cost is a one-time, compile-time expense that no runtime tuning loop contributes to, confirming that \name's analytical formulation scales to large, real-world kernels and workloads.
